\subsection{Матричная форма системы уравнений движения в изображениях по Лапласу}

Применим оператор Лапласа к системе уравнений движения при нулевых начальных условиях:
\begin{equation}
\begin{cases}
    \left(
        J_{\text{дпр}}p^2
        +
        bp
        +
        c
    \right)
    \Phi_{\text{рпр}}
    -
    \left(
        bp+c
    \right)
    \Phi_{\text{м}}
    -
    M_{\text{дпр}}
    =
    0,
    \\[0.4em]

    \left(
        J_{\text{м}}p^2
        +
        bp
        +
        c
    \right)
    \Phi_{\text{м}}
    -
    \left(
        bp+c
    \right)
    \Phi_{\text{рпр}}
    =
    0,
    \\[0.4em]

    \left(
        \tau p+1
    \right)
    M_{\text{дпр}}
    +
    s_{\text{пр}}p
    \Phi_{\text{рпр}}
    -
    r_{\text{пр}}U
    =
    0,
    \\[0.4em]

    pU
    +
    kk_{\text{п}}p
    \Phi_{\text{рпр}}
    +
    kk_{\text{и}}
    \Phi_{\text{м}}
    =
    kk_{\text{и}}
    \Phi^*.
\end{cases}
\end{equation}

Представим систему в матричной форме:
\[
    A \cdot X = h \Phi^*.
\]

Тогда матрица системы имеет вид
\begin{equation}
    A=
    \begin{bmatrix}
        J_{\text{дпр}}p^2+bp+c
        &
        -(bp+c)
        &
        -1
        &
        0
        \\[0.4em]

        -(bp+c)
        &
        J_{\text{м}}p^2+bp+c
        &
        0
        &
        0
        \\[0.4em]

        s_{\text{пр}}p
        &
        0
        &
        \tau p+1
        &
        -r_{\text{пр}}
        \\[0.4em]

        kk_{\text{п}}p
        &
        kk_{\text{и}}
        &
        0
        &
        p
    \end{bmatrix}.
\end{equation}

Вектор неизвестных:
\[
    X=
    \begin{bmatrix}
        \Phi_{\text{рпр}}
        \\
        \Phi_{\text{м}}
        \\
        M_{\text{дпр}}
        \\
        U
    \end{bmatrix}.
\]

Вектор внешнего воздействия:
\[
    h=
    \begin{bmatrix}
        0
        \\
        0
        \\
        0
        \\
        kk_{\text{и}}
    \end{bmatrix}.
\]
